\documentclass[11pt]{article}
\usepackage[margin=0.85in]{geometry}
\usepackage{amsmath,booktabs,xurl,hyperref}
\usepackage[T1]{fontenc}
\setlength{\emergencystretch}{2em}
\title{Encounter Range and a Terminal Departure Certificate}
\author{Pan Dynamics Collision Avoidance Research}
\date{October 1, 2026}
\begin{document}
\maketitle

\section{Summary}
The target motion model (constant speed rate and constant sideslip) is now
assumed only during an encounter. An encounter ends when the ego and target
bodies are more than \path{cfg.collision.encounterRangeMeters} apart; the
default is 30 m, a declared modeling threshold. The controller no longer
requires separation from the complete future target motion as its ordinary
terminal condition. An endpoint that is still inside the range must instead
show that its own terminal policy leaves the range while separated from the
target at matching times.

At 8 m/s the first-frame prediction of \texttt{turningCrossing} falls from 216
to 93 holds and that of \texttt{curvedCrossing} from 254 to 90. All fourteen
closed-loop runs (seven scenarios at 8 and 15 m/s) complete collision-free
with zero hard residual and zero safety slack and are confirmed recovered.
Nominal recovery is faster in three runs, unchanged in seven and slower in
four. One of the four, \texttt{curvedCrossing} at 15 m/s, is markedly worse
and is diagnosed in Section~\ref{sec:regression}. The longest controller frame
is unchanged at 5.1--5.4 s: it is set by the five-second improvement budget,
not by the prediction length.

\section{Design}
\subsection{Encounter exit inside the prediction}
Let $d_j$ be the rectangle distance between the ego body at prediction node
$j$ and the target body at the same absolute time, and $R$ the range. The
exit is the first node with $d_j>R$. Holds up to the exit keep every collision
row, interior node and interval certificate; later holds carry no target
constraint. A target that is out of range at the initial node imposes no
constraint in that call; every call examines its own initial node again, and
moving-flow seeds are screened at the first engaged sample. A retained
prediction shifts its exit with it.

\subsection{Departure of the endpoint policy}
When the prediction is still inside the range at its endpoint, the horizon is
not lengthened. The reference pose of the fitted terminal core is followed at
matching absolute times until the rectangle distance, less the enclosed core
deviation $\epsilon$, exceeds $R$. With $\delta=d_{\min}+2\epsilon$ the
certificate requires $d_j>\delta$ at every visited node. Nodes are visited in
strides: over a time $s$ the distance changes by at most
$(v_g+(|V_T|+|A_T|s)(1+|\kappa_T|r_T))s$, and a stride is the largest whole
number of holds for which this stays below $d_j-\delta$, so the reference
remains separated by $\delta$ between visited nodes. A node closer than one
such hold uses the existing interval certificate. The closest visited node
supplies four polygon-dual rows for the conic and correction programs. A
retained pose reuses its certified exit node and margin; a refitted pose is
followed again.

If the range is never exceeded (a target travelling alongside), the earlier
sufficient all-future separation remains the terminal requirement. It is
also accepted whenever it holds. \texttt{encounterRangeMeters=Inf} never ends
an encounter and recovers the previous requirement.

\subsection{A discarded first variant}
The first implementation required the exit to occur inside the inherited
prediction. Because the prediction length is inherited and only shortens, the
optimizer was forced to be 30 m from the target by a fixed absolute time. In
\texttt{turningCrossing} at 8 m/s this produced a 10-m lateral detour with
acceleration to 12 m/s, a maximum CLF slack of 474 (49 before) and a
17.35-s recovery (16.0 s before). Its seven 8-m/s results are retained in
\path{closed-loop.csv} as \texttt{exitInsideHorizon}; it was replaced by the
departure certificate above.

\section{First-frame prediction length}
\begin{center}
\begin{tabular}{lrrrr}
\toprule
 & \multicolumn{2}{c}{8 m/s} & \multicolumn{2}{c}{15 m/s} \\
Scenario & $R=\infty$ & $R=30$ m & $R=\infty$ & $R=30$ m \\
\midrule
headOn & 92 & 92 & 96 & 96 \\
acceleratingHeadOn & 92 & 92 & 95 & 95 \\
brakingLead & 195 & 195 & 196 & 196 \\
crossing & 93 & 93 & 96 & 96 \\
turningCrossing & 216 & 93 & 136 & 93 \\
curvedHeadOn & 123 & 92 & 98 & 90 \\
curvedCrossing & 254 & 90 & 156 & 96 \\
\bottomrule
\end{tabular}
\end{center}
Holds are 50 ms. \texttt{brakingLead} is unchanged: its length comes from the
moving-flow guidance staying active beside the lead vehicle, not from the
terminal target requirement. In \texttt{curvedHeadOn} at 15 m/s the target
starts beyond 30 m and imposes no constraint in the first frame.

\section{Closed-loop campaign}
Each run uses \texttt{runNonlinearPredictiveSafetyValidation} with the original
fixtures, a tight ODE45 plant replay, no road constraints, observer, noise or
random draw, and stops after a five-second confirmed nominal recovery.
``Before'' is the campaign of \path{CLF_NOMINAL_RECOVERY_20261001.tex}.
\begin{center}
\small
\begin{tabular}{llrrrrrr}
\toprule
 & & \multicolumn{2}{c}{Recovery (s)} & \multicolumn{2}{c}{Controller (s)} & \multicolumn{2}{c}{Max CLF slack} \\
Speed & Scenario & before & after & before & after & before & after \\
\midrule
8 & headOn & 13.05 & 13.05 & 171.9 & 181.1 & 107.8 & 171.5 \\
8 & acceleratingHeadOn & 13.05 & 13.05 & 182.5 & 175.6 & 175.5 & 180.1 \\
8 & brakingLead & 14.80 & 15.10 & 466.8 & 426.4 & 66.0 & 71.1 \\
8 & crossing & 13.05 & 13.05 & 176.7 & 177.4 & 66.3 & 34.4 \\
8 & turningCrossing & 16.00 & 13.05 & 213.1 & 203.9 & 49.4 & 51.8 \\
8 & curvedHeadOn & 19.60 & 13.05 & 304.7 & 207.9 & 188.1 & 181.2 \\
8 & curvedCrossing & 17.90 & 19.00 & 326.9 & 368.0 & 34.5 & 64.0 \\
15 & headOn & 13.05 & 13.05 & 140.4 & 112.9 & 81.2 & 172.9 \\
15 & acceleratingHeadOn & 14.75 & 15.55 & 146.0 & 139.1 & 166.1 & 141.1 \\
15 & brakingLead & 16.15 & 16.15 & 505.9 & 505.3 & 49.0 & 49.0 \\
15 & crossing & 13.05 & 13.05 & 197.4 & 172.4 & 73.6 & 104.6 \\
15 & turningCrossing & 15.60 & 13.05 & 226.7 & 206.4 & 74.4 & 108.5 \\
15 & curvedHeadOn & 13.05 & 13.05 & 200.3 & 167.3 & 74.4 & 86.2 \\
15 & curvedCrossing & 16.45 & 18.45 & 238.7 & 655.4 & 71.7 & 1711.0 \\
\bottomrule
\end{tabular}
\end{center}
13.05 s is the earliest possible confirmation (eight-second minimum plus the
five-second dwell). Every run has zero maximum hard residual and zero maximum
safety slack. The minimum replay clearance over the fourteen after-runs is
0.0060 m against the required 0.006 m (\texttt{brakingLead} at both speeds and
\texttt{curvedCrossing} at 15 m/s sit at the margin); before, it was 0.0060 m.
Controller time is a sum over frames. The before-campaign ran overlapping
processes and the after-campaign ran one process at a time; frames that reach
the five-second budget depend on machine load, so these times are functional
measurements, not an isolated benchmark.

\section{The slower curved crossing}\label{sec:regression}
In \texttt{curvedCrossing} at 15 m/s the ego brakes to 2.4 m/s, turns after
the target and returns later; it stays collision-free, but the maximum CLF
slack is 1711 and 129 frames take at least one second (42 before). Three
additional runs of the same scenario isolate the cause (recovery in s, speed in m/s):
\begin{center}
\begin{tabular}{lrrrrr}
\toprule
Run & Recovery & CLF slack & Min speed & Slow frames & Min holds \\
\midrule
before & 16.45 & 71.7 & 10.5 & 42 & 76 \\
this change, $R=30$ m & 18.45 & 1711.0 & 2.4 & 129 & 76 \\
this change, $R=\infty$ & 16.55 & 75.2 & 9.7 & 48 & 76 \\
this change, $R=30$ m, 6-s completion & 14.15 & 57.5 & 10.1 & 37 & 136 \\
\bottomrule
\end{tabular}
\end{center}
With $R=\infty$ the new code reproduces the earlier behavior, so the
difference is caused by the range, not by the other edits. With the range and
a six-second completion allowance (\texttt{recoveryHorizonSeconds=6}) the run
is better than before in every listed measure. The earlier conservative
certificate had an unintended side effect: its long seed (156 holds, 7.8 s)
gave the optimizer a long look-ahead during the encounter. With the departure
certificate the prediction is at its configured minimum of 76 holds (3.8 s)
after 1.2 s, and the endpoint must already be inside the small straight-cruise
core at 15 m/s by then. That is too short for this maneuver. The same
mechanism plausibly explains the 1.1-s slower recovery of the 8-m/s curved
crossing, which was not separately tested.

The default \texttt{recoveryHorizonSeconds} is left at 3 s in this change. A
longer minimum prediction raises the cost of every frame, and only one
scenario was run with 6 s; the default should be chosen from a full campaign.

\section{Validation}
Thirteen controller test files, 313 cases, pass on the final source
(\path{tests.csv}). New cases cover: a target beyond the range imposes no
constraint; a target entering the range is avoided from that sample using a
moving-flow seed; an in-horizon exit shifts with a retained prediction; an
unbounded range keeps the all-future requirement; the departure margin for a
receding, a blocking and an alongside target; and, for the turning crossing,
an independent ODE45 rollout of the terminal policy from the endpoint that
stays separated until the range is exceeded. The 8-m/s campaign ran on a
source differing from the final one by one guard on a missing witness field;
the saved traces show that branch is not reached in any campaign frame. The
15-m/s campaign and the diagnostic runs used the final source.

\section{Limits}
The range is an assumption about when a target matters, not a derived bound.
After an exit the certificate says nothing about the target; safety against a
target that returns rests on the later call in which it is again within
range, and the shift argument does not cover the first sample of a new
encounter. The exit is detected on the hold grid and at visited departure
nodes. The departure margin is the closest visited distance and is not a
smooth function of the endpoint when the visited nodes change. All results
are nominal-model predictions with a sampled ODE45 replay; they are not a
continuous-plant certificate. Estimator-in-the-loop scenarios and the
PassVeh14DOF scenarios were not run.
\end{document}
